
\subsubsection{6.1 Interpretation of the Efficiency Advantage}

The efficiency ratio of approximately $47\times$ (central estimate) or at least 6.$8\times$ (conservative lower bound) indicates that structural permutation equivariance substantially reduces the number of training models required for weight-space property prediction. GNN-NFN reaches 90\% of its peak $R^2$ at approximately $N\approx 147$ models; flat-MLP does not reach 90\% of its own peak $R^{2}=0.886$ within any sampled training size up to N=1,000, requiring approximately the full training set (~7,000 models).

This finding has a practical interpretation: an organization collecting model checkpoints for weight-space analysis needs approximately 250 models to obtain strong performance from GNN-NFN ($R^{2}=0.767)$, compared to approximately 7,000 models for flat-MLP to reach comparable relative performance. Whether this gap reflects a structural hypothesis-space reduction (equivariant architectures search within permutation-symmetric function classes) or gradient dynamics (flat-MLP eventually discovers symmetric solutions but requires more data to do so) cannot be determined from present evidence. The causal mechanism --- Step 2 of the chain linking structural equivariance to efficiency --- was not directly verified. The 6.$8\times$ ratio is consistent with the hypothesis-space reduction interpretation but does not prove it.

\subsubsection{6.2 Interpretation of the Data-Regime Crossover}

At N=100, GNN-NFN achieves $R^{2}=-0.016$, performing below the chance level of predicting the mean. The structural inductive bias, expected to help most when data is scarce, is instead harmful at this scale. Two explanations are plausible:

\textbf{Underfitting due to architecture complexity (plausibility: high).} GNN-NFN's graph neural network architecture requires a sufficient diversity of weight-graph topologies in its training set to learn useful node embeddings. At N=100, insufficient topological diversity exists for the encoder to generalize. The sharp jump from $R^{2}=-0.016$ (N=100) to $R^{2}=0.767$ (N=250) is consistent with crossing a minimum-data threshold for graph encoder generalization, not with the smooth degradation expected from learning-rate instability.

\textbf{Hyperparameter mismatch at small N (plausibility: medium).} Adam optimizer hyperparameters were tuned on the full-data condition. At N=100, the learning rate may cause the architecturally more complex GNN-NFN to diverge or fail to converge, while the simpler flat-MLP is more robust to suboptimal hyperparameters at small batch counts.

An LR sweep for GNN-NFN at N=100 is the highest-priority ablation to distinguish these explanations. If GNN-NFN achieves positive $R^2$ under a lower learning rate at N=100, the hyperparameter explanation is supported; if it remains near $R^{2}=0$ across learning rates, the underfitting explanation is more likely.

PermAug's advantage at N=100 appears primarily due to the $11\times$ data expansion: 100 base models become 1,100 effective training samples. The PermAug expansion factor ($11\times )$ is a fixed design choice that was not ablated; a reviewer may ask whether the crossover at N=100 depends on this expansion factor or whether it is driven by the structural difference between augmentation and equivariance. This ablation is a high-priority follow-up.

\subsubsection{6.3 Connection to Expressivity Theory}

Full-scale convergence ($\Delta \leq 0.038 R^{2})$ provides empirical support for the expressivity equivalence theorem of Dayan et al. [2026] and characterizes the practical window of equivariant advantage: approximately N $\in$ [150, 7,000] on the CIFAR-10 CNN zoo. Below this window, structural equivariance provides no advantage (and may be harmful); above the upper end, both encoders approach similar ceiling performance. The theoretical result establishes the ceiling parity; the present data characterize the efficiency window between the two endpoints.

\subsubsection{6.4 Distinguishability of Augmentation and Structural Equivariance}

The N=100 crossover refutes the original Assumption A3 that permutation augmentation and structural equivariance are interchangeable approximations. At N=100, they produce qualitatively different outcomes: PermAug is beneficial ($R^{2}=0.138)$ while structural equivariance is harmful ($R^{2}=-0.016)$. This is a positive finding: the two symmetry-enforcement strategies are empirically distinguishable, and their relative performance depends on data regime. Practitioners with very small model zoos (<150 models) should prefer augmentation of a plain encoder; those with larger zoos ($\geq 250$ models) should prefer structural equivariance.

\subsubsection{6.5 Limitations}

\textbf{CIFAR-10 only.} All property-prediction efficiency results are from the CIFAR-10 CNN zoo. The MNIST MLP zoo --- for which DWSNets is architecturally suited --- was not available locally and was not evaluated. The efficiency ratio of 6.$8\times--47\times$ applies specifically to the CIFAR-10 CNN zoo setting. Whether the same ratio holds for MLP zoos (where DWSNets is the appropriate equivariant encoder) is unknown.

\textbf{DWSNets excluded from property prediction.} DWSNets requires M>2 FC layers; the CIFAR-10 CNN zoo has 2 FC layers. The property-prediction efficiency advantage was measured for GNN-NFN only. DWSNets equivariance is structurally confirmed on synthetic MLP weights, but DWSNets property-prediction efficiency on CNN zoos is not reported.

\textbf{N=100 crossover is single-seed.} The most novel finding lacks multi-seed confidence intervals. The N=100 values for PermAug and GNN-NFN are point estimates with collapsed CI\_lo = CI\_hi. This finding is presented as a preliminary observation, not a confirmed result. H-M3 is flagged as PARTIAL/LIMITATION\_RECORDED in the research pipeline.

\textbf{PermAug from separate experimental phase.} PermAug results come from H-M3, a separate experiment that loaded baseline flat-MLP and GNN-NFN results from H-M2. PermAug training used a distinct random seed. The crossover finding should be confirmed under a fully co-trained single-experiment design.

\textbf{Low zoo diversity.} CIFAR-10 zoo accuracy variance=0.025. Models cluster near convergence, potentially making the prediction task easier for all encoders and inflating absolute $R^2$ values. The relative efficiency ratio is preserved but absolute $R^2$ values should not be generalized to harder, more diverse zoos.

\textbf{Hyperparameter fairness at small N.} Adam hyperparameters were tuned at full data and held fixed. At N=100, this may disadvantage GNN-NFN's more complex architecture. An LR sweep at N=100 is needed.

\textbf{No formal CI on efficiency ratio.} The interpolated N\_equiv,$90\approx 147$ spans the interval [101, 249], yielding efficiency ratios from approximately $28\times$ to $69\times$ using N\_plain,90=7,000. The central estimate of $47\times$ is a point estimate under linear interpolation; a formal bootstrap CI on the efficiency ratio requires multi-seed test-set resampling.

\textbf{PermAug expansion factor not ablated.} The $11\times$ expansion factor is a fixed design choice. Whether the N=100 crossover depends on this factor --- or would occur at a lower expansion --- is not determined.

